\chapter{Родитель собственного четырёхмерного объёма ячейки}
\label{ch:v10-cell-birth-intrinsic-four-volume-parent-origin}

\section{Геометрический носитель}

После отрицательного аудита требуется новый внутренний размерный объект.
Минимальная изотропная четырёхмерная ячейка задаётся грамовой матрицей
\begin{equation}
 G_{\rm cell}=\ell_{\rm cell}^2I_4,
 \qquad \det G_{\rm cell}=\ell_{\rm cell}^8,
 \qquad v_{\rm cell}=\sqrt{\det G_{\rm cell}}=\ell_{\rm cell}^4.
 \label{eq:v10-cell-four-volume}
\end{equation}
Нормированная форма
\begin{equation}
 R_{\rm cell}=\frac{G_{\rm cell}}{\sqrt{v_{\rm cell}}}=I_4,
 \qquad \det R_{\rm cell}=1
 \label{eq:v10-cell-unimodular-shape}
\end{equation}
отделяет форму от размера. При рождении одной ячейки
\begin{equation}
 V_N=Nv_{\rm cell},
 \qquad V_{N+1}-V_N=v_{\rm cell}.
 \label{eq:v10-cell-volume-birth-law}
\end{equation}
Таким образом, собственный объём является корректным локальным носителем
приращения, а не повторным обозначением полного расширения.

\section{Безразмерный объёмно-энергетический инвариант}

Из объёма и энергии часов строится
\begin{equation}
 Y=\frac{E_C^4v_{\rm cell}}{(\hbar c)^4}.
 \label{eq:v10-cell-volume-energy-invariant}
\end{equation}
Условие $Y=1$ эквивалентно комптоновской связи
\begin{equation}
 E_C\ell_{\rm cell}=\hbar c.
 \label{eq:v10-cell-compton-relation}
\end{equation}
Однако преобразование
\begin{equation}
 \ell_{\rm cell}\mapsto s\ell_{\rm cell},
 \qquad E_C\mapsto\frac{E_C}{s}
 \label{eq:v10-cell-volume-energy-orbit}
\end{equation}
сохраняет $Y$. Следовательно, даже точное происхождение отношения $Y=1$
не выбирает величины двух множителей по отдельности.

\section{Общий родитель и его ядро}

Пусть $u=\chi^2$, $\rho=\Gamma_B/\Omega$, а $\epsilon$ и $\lambda$
являются логарифмическими вариациями энергии и длины. Совместный родитель
\begin{equation}
 \mathcal P_{vE}
 =\frac12(u-k_X)^2+\frac12(\rho-u)^2
 +\frac12(\epsilon+\lambda)^2
 \label{eq:v10-cell-volume-common-parent}
\end{equation}
имеет гессиан
\begin{equation}
 \operatorname{Hess}\mathcal P_{vE}=
 \begin{pmatrix}
 2&-1&0&0\\-1&1&0&0\\0&0&1&1\\0&0&1&1
 \end{pmatrix}.
 \label{eq:v10-cell-volume-parent-hessian}
\end{equation}
Его ранг равен $3$, ядро одномерно и порождается вектором
$(0,0,1,-1)^T$. Спектр равен
\begin{equation}
 \left\{0,2,\frac{3-\sqrt5}{2},\frac{3+\sqrt5}{2}\right\}.
 \label{eq:v10-cell-volume-parent-spectrum}
\end{equation}
Родитель положительно полуопределён и выбирает все относительные данные,
но оставляет точно одну масштабную моду.

\begin{theorem}[Собственный объём как относительный геометрический носитель]
Четырёхмерная ячейка имеет точное приращение объёма $v_{\rm cell}$, а
общий родитель совместно фиксирует прежнюю связь рождения и безразмерное
произведение $E_C^4v_{\rm cell}/(\hbar c)^4$. Архитектура замыкается, но
гессиан сохраняет одно масштабное ядро.
\end{theorem}

\begin{nogocriterion}[Комптоновская связь не выбирает размер ячейки]
Равенство $E_C\ell_{\rm cell}=\hbar c$ определяет только произведение.
Без независимого спектрального, топологического или мерного правила,
выбирающего $v_{\rm cell}$, оно не выводит ни абсолютную длину, ни энергию
часов.
\end{nogocriterion}

\begin{protocol}[Статус собственного объёма]\begin{enumerate}
\item четырёхмерный метрический носитель: \textbf{построен};
\item закон приращения объёма при рождении: \textbf{получен};
\item архитектура: \textbf{$8/8$};
\item относительное происхождение: \textbf{$3/3$};
\item ранг и размерность ядра общего гессиана: \textbf{$3/1$};
\item величина собственного объёма: \textbf{$0/1$};
\item абсолютная энергия часов: \textbf{$0/1$};
\item следующий гейт: \textbf{спектральная счётная мера четырёхмерного
объёма ячейки}.
\end{enumerate}\end{protocol}

Машинный сертификат сохранён в
\path{s2t_v10_cell_birth_intrinsic_four_volume_parent_origin_gate_results.json}.